\providecommand{\kBKM}{\kappa_{\mathrm{BKM}}}
\providecommand{\kch}{\kappa_{\mathrm{ch}}}
\providecommand{\kcat}{\kappa_{\mathrm{cat}}}
\providecommand{\kfib}{\kappa_{\mathrm{fiber}}}
\providecommand{\Sp}{\mathrm{Sp}}
\providecommand{\GSpin}{\mathrm{GSpin}}
\providecommand{\Sh}{\mathrm{Sh}}
\providecommand{\KM}{\mathrm{KM}}
\providecommand{\wt}{\mathrm{wt}}
\providecommand{\Hol}{\mathrm{Hol}}
\providecommand{\PP}{\mathrm{PP}}
\providecommand{\cO}{\mathcal{O}}
\providecommand{\bH}{\mathbb{H}}
\providecommand{\bQ}{\mathbb{Q}}
\providecommand{\bZ}{\mathbb{Z}}
\providecommand{\DeltaFive}{\Delta_5}
\providecommand{\Dfive}{D_5}

\section*{Kudla--Millson Comparison at CHL Levels}
\label{sec:kudla-millson-chl-levels}

\subsection*{Three separate inputs}

The primitive CHL denominator ladder is indexed by
\[
  N\in\{1,2,3,4,6\}.
\]
Its Borcherds input is the \(g_N\)-twined weak Jacobi form
\(\phi^{(N)}_{0,1}\) of weight \(0\) and index \(1\).  In the
standard CHL product normalisation its discriminant-zero constants are
\[
  (c_1(0),c_2(0),c_3(0),c_4(0),c_6(0))
  =
  (10,8,6,4,2),
\]
and Borcherds' singular-theta theorem gives
\[
  \wt(\Phi_N)=\kBKM(\Phi_N)=\frac{c_N(0)}2
  =
  (5,4,3,2,1)
\]
for \(N=(1,2,3,4,6)\).  The level-one product is
\(\DeltaFive\), and the monic denominator convention is
\[
  \Dfive=64^{-1}\DeltaFive .
\]

The Gritsenko--Cl\'ery diagonal-divisor catalogue is a different
family.  It is indexed by paramodular triples
\[
 (t,N_{\mathrm{dd}};k)=
 (1,1;5),(2,1;2),(1,2;3),(3,1;1),(1,3;2),
 (4,1;\tfrac12),(1,4;\tfrac32),(2,2;1),
\]
with twice-weight tuple
\[
  (10,4,6,2,4,1,3,2).
\]
The common entry is the row \((1,1;5)\), namely \(\DeltaFive\).  At the
next paramodular row the weight is \(2\), while the CHL \(N=2\)
denominator has weight \(4\).  These two numbers already exclude an
identity between the full CHL ladder and the diagonal-divisor catalogue.
Any further comparison requires a row-map theorem identifying the
automorphic group, cover, character, divisor, Jacobi input, Weyl
chamber, and leading coefficient.

The Kudla--Millson construction is a third datum.  It supplies
differential forms and theta kernels on orthogonal Shimura varieties.
It computes special-cycle cohomology classes and regularised theta
lifts after a lattice, Weil representation, and principal part have
been fixed.  Products with different weights remain distinct.

\subsection*{Correct theta-lift statement}

Fix \(N\in\{2,3,4,6\}\).  Let \(L_N\) be an even lattice of signature
\((2,n_N)\) with Weil representation \(\rho_{L_N}\), and let
\[
  f_N\in M^{!}_{1-n_N/2}(\rho_{L_N})
\]
be a weakly holomorphic vector-valued modular form satisfying the
integrality and principal-part hypotheses of Borcherds' product
theorem.  If the principal part of \(f_N\) is the CHL principal part
whose Borcherds product is \(\Phi_N\), then the regularised theta lift
\[
  \Psi(f_N)(Z)
  =
  \exp\!\left(
  \int_{\Gamma\backslash\bH}^{\mathrm{reg}}
  \langle f_N(\tau),\Theta_{L_N}(\tau,Z)\rangle
  \frac{dx\,dy}{y^2}
  \right)
\]
is a meromorphic automorphic product with divisor determined by that
principal part and with weight
\[
  \wt(\Psi(f_N))=\frac{c_N(0)}2.
\]
When \(\Psi(f_N)\) is normalised to the CHL denominator, it equals
\(\Phi_N\).  This is the precise content supplied by Borcherds'
regularised product theory and Bruinier's reconstruction of products
from principal parts.

The Kudla--Millson form \(\Theta_{\KM,L_N}\) represents the same
special-cycle class in cohomology.  After holomorphic projection, a
scalar-valued comparison with \(\Phi_N\) has the form
\[
  \Hol\!\left[
  \int_{\Gamma\backslash\bH}^{\mathrm{reg}}
  \langle f_N(\tau),\Theta_{\KM,L_N}(\tau,Z)\rangle
  \frac{dx\,dy}{y^2}
  \right]
  =
  \lambda_N\,\Phi_N
\]
only under the following hypotheses:
\[
\begin{array}{ll}
\text{(i)}&\text{the lattice }L_N\text{ and the CHL product lattice agree;}\\
\text{(ii)}&\text{the principal part is the CHL principal part;}\\
\text{(iii)}&\text{the target cusp-form space is one-dimensional;}\\
\text{(iv)}&\text{the leading Fourier coefficient fixes }\lambda_N=1;\\
\text{(v)}&\text{the character and chamber match the CHL denominator.}
\end{array}
\]
These five checks are the condition for upgrading the formula from a
theta-lift comparison problem to a theorem identifying a named
paramodular form.

\subsection*{Level-one normalisation}

At \(N=1\), the Eichler--Zagier generator satisfies
\[
  \phi_{0,1}^{\mathrm{EZ}}(\tau,z)|_{q^0}
  =
  \zeta^{-1}+10+\zeta .
\]
Hence \(c_1(0)=10\), and
\[
  \kBKM(\DeltaFive)=\wt(\DeltaFive)=5.
\]
The geometric K3 elliptic genus is \(2\,\phi_{0,1}^{\mathrm{EZ}}\).  Its
Borcherds lift gives the Igusa square
\(\Phi_{10}^{\mathrm{un}}=\DeltaFive^2\), while the
Oberdieck--Pandharipande scalar convention uses
\[
  \Phi_{10}^{\mathrm{OP}}=\Dfive^2=4096^{-1}\DeltaFive^2 .
\]
Thus the primitive denominator, the Igusa square, and the reduced scalar
partition function have different normalisations.

\subsection*{Compact \(K3\times E\) invariants}

For \(Y=K3\times E\), Kunneth gives
\[
  h^{0,\bullet}(Y)=(1,1,1,1),\qquad
  \kch(Y)=1-1+1-1=0.
\]
The categorical invariant is
\[
  \kcat(Y)=\chi(\cO_{K3})\chi(\cO_E)=2\cdot0=0.
\]
The remaining distinguished readings are
\[
  \kch^{\mathrm{Heis}}(Y)=3,\qquad
  \kBKM(\DeltaFive)=5,\qquad
  \kfib(K3)=24.
\]
These values come from the compact Hodge supertrace, the total-space
Euler characteristic of the structure sheaf, the Heisenberg--Mukai
specialisation, the Borcherds denominator weight, and the K3 Mukai
fibre rank.

\subsection*{Proof obligations at \(N=2,3,4,6\)}

The CHL products \(\Phi_N\) have verified weights
\[
  \kBKM(\Phi_2)=4,\qquad
  \kBKM(\Phi_3)=3,\qquad
  \kBKM(\Phi_4)=2,\qquad
  \kBKM(\Phi_6)=1.
\]
A rank-three Borcherds--Kac--Moody superalgebra theorem at these levels
requires more than the automorphic product.  One must construct
\[
  (L_N,\Delta_N^{\mathrm{re}},\rho_N,W_N,\epsilon_N,m_N,\Phi_N),
\]
where \(L_N\) is the Lorentzian root lattice, \(\Delta_N^{\mathrm{re}}\)
is the set of real simple roots, \(\rho_N\) is the Weyl vector, \(W_N\)
is the reflection group, \(\epsilon_N\) is the reflection character,
and \(m_N\) is the integral signed root-multiplicity function.  The
denominator identity must then be proved with the same product-degree
map and chamber used in the automorphic construction.

The Kudla--Millson comparison supplies a cohomological route to the
special divisor.  It leaves the root datum, parity lift, Lie bracket,
and denominator identity as separate algebraic constructions.

\paragraph{References.}
Borcherds, \emph{Automorphic forms with singularities on
Grassmannians}, Invent.\ Math.\ 132 (1998), Theorem~13.3.  Bruinier,
\emph{Borcherds Products on O(2,l) and Chern Classes of Heegner
Divisors}, Lecture Notes in Mathematics 1780 (2002).  Kudla--Millson,
\emph{Intersection numbers of cycles on locally symmetric spaces and
Fourier coefficients of holomorphic modular forms in several complex
variables}, Publ.\ IHES 71 (1990).  Gritsenko--Nikulin 1995 and 1998
for the \(\DeltaFive\) denominator.  Gritsenko 1999 for the arithmetic
lift normalisation.  Gritsenko--Cl\'ery 2013 for the eight
diagonal-divisor forms.
